% Copyright (c) 2026 Geovana Neves. Licensed under CC BY 4.0 (https://creativecommons.org/licenses/by/4.0/). See LICENSE-AND-AUTHORSHIP.md.
%
% 03-gui-to-pyfs/text/05-flight-and-solver.tex
%
% Transcribed from the tables of docs/gui-to-pyfs.md, one row per GUI step:
% the key that takes it (links and solver commands dropped) and the builds
% its commands are verified on. When the page changes a row, change the same
% row here, and the counts of 01-reading.tex with it.

\section{Flight conditions and solver}

\begin{frame}{Flight conditions and solver: what a key takes (1/6)}
\relax
{\ftstablesize\renewcommand{\arraystretch}{1.05}
\begin{tabularx}{\linewidth}{@{}JKZ@{}}
\guihead
Set the free stream: its speed, angle of attack and sideslip & row keys \code{ALPHA} and \code{BETA} in the \code{FLIGHT\_CONDITION} cell, whose \code{MACH} or \code{TASmps} gives the speed; row key \code{VELOCITY} from Python alone & 26.120 to 26.124, except one \\
Set a custom free stream: import a velocity field over the YZ plane from a file & row key \code{FREESTREAM}, the stem of a file of \code{inputs/freestreams/}, a \code{.txt} in the STRUCTURED form or a \code{.dat} in the UNSTRUCTURED form, in m and m/s; the field is the flow's direction, so the row's angle of attack and sideslip are 0 & none \\
Sweep the angle of attack or the sideslip & \code{ALPHA:sweep} or \code{BETA:sweep} in the \code{FLIGHT\_CONDITION} cell, with the \code{SWEEP\_VALUES} column: one script; every point starts cold, clearing the solution before it, and row key \code{COLD\_START} set to \code{false} keeps the earlier warm start from the last point & 26.120 to 26.124 \\
Set the fluid: density, pressure, temperature, viscosity, speed of sound & the \code{FLIGHT\_CONDITION} cell, or the setup's \code{[flight\_condition]} table & 26.101 to 26.124 \\
\bottomrule
\end{tabularx}}
\guisource{Flight conditions and solver}
\end{frame}

\begin{frame}{Flight conditions and solver: what a key takes (2/6)}
\relax
{\ftstablesize\renewcommand{\arraystretch}{1.05}
\begin{tabularx}{\linewidth}{@{}JKZ@{}}
\guihead
Set the reference area, length and velocity & reference lengths; setup key \code{reference\_velocity\_m\_per\_s} & 26.101 to 26.124 \\
Give the free stream as a Mach number, set the reference Mach number or the speed of sound, or switch the reference velocity off & setup keys \code{freestream\_input}, \code{reference\_mach} and \code{disable\_reference\_velocity}; setup key \code{sonic\_velocity\_m\_per\_s} is refused on every build, since no build records \code{SONIC\_VELOCITY} and the sound speed follows from temperature and specific-heat ratio & 26.120 to 26.124, except two \\
State the units of a custom free-stream file & row key \code{FREESTREAM\_UNITS}: \code{SI} for metres and metres per second, which the package scales once into the simulation's unit, or \code{NATIVE} for a file already in it & none \\
Choose steady or unsteady, the time step and the number of steps & the \code{WORKFLOW} column: \code{steady}, \code{unsteady}, \code{unsteady\_rotor} or \code{qsteady\_rotor}; row keys \code{DELTA\_TIME} and \code{TIME\_ITERATIONS}, or \code{DELTA\_THETA} and \code{REVOLUTIONS} on a rotor row & 26.100 to 26.124 \\
Set the iterations and the convergence threshold & setup keys \code{iterations}, \code{convergence}, \code{forced\_iterations} and \code{convergence\_iterations} & 26.120 to 26.124 \\
\bottomrule
\end{tabularx}}
\guisource{Flight conditions and solver}
\end{frame}

\begin{frame}{Flight conditions and solver: what a key takes (3/6)}
\relax
{\ftstablesize\renewcommand{\arraystretch}{1.05}
\begin{tabularx}{\linewidth}{@{}JKZ@{}}
\guihead
Set the number of processors & row key \code{NCPUS}, or setup key \code{max\_threads} & 26.120 to 26.124 \\
Choose the boundary layer and its coupling & setup keys \code{boundary\_layer} and \code{viscous\_coupling} & 26.120 to 26.124 \\
Switch the viscous coupling on at a time step of an unsteady run & setup key \code{unsteady\_viscous\_coupling\_iteration}, refused on 26.124, which does not answer the command & none \\
Initialise the solver: the compressibility model, wall-collision avoidance, a mirror plane or periodic copies & setup keys \code{solver\_model} and \code{wall\_collision\_avoidance}; row keys \code{SYMMETRY} and \code{PERIODIC\_COPIES}; setup key \code{legacy\_solver\_model}, the older model command, refused on every build: only 25.000 documents \code{SET\_SOLVER\_MODEL} and no workflow writes that edition's \code{INITIALIZE\_SOLVER} & 26.101 to 26.124, except one \\
Remove a saved initialization, mark proximal boundaries, set the automatic trailing-edge and wake-node physics, or delete transition trips & setup keys \code{remove\_initialization}, \code{proximal\_boundaries}, \code{physics\_auto\_trailing\_edges}, \code{physics\_auto\_wake\_nodes} and \code{delete\_transition\_trips} & 26.121 to 26.124, except two \\
\bottomrule
\end{tabularx}}
\guisource{Flight conditions and solver}
\end{frame}

\begin{frame}{Flight conditions and solver: what a key takes (4/6)}
\relax
{\ftstablesize\renewcommand{\arraystretch}{1.05}
\begin{tabularx}{\linewidth}{@{}JKZ@{}}
\guihead
Advanced settings: far-field layers, the mesh-induced wake velocity, wake-on-wake induction, a further wake relaxation & setup keys \code{farfield\_layers}, \code{mesh\_induced\_wake\_velocity}, \code{wake\_on\_wake\_induction} and \code{additional\_wake\_relaxation} & 26.120, 26.122, 26.123 \\
Advanced settings: the minimum pressure coefficient, stabilisation, unsteady pressure and Kutta condition, Reynolds-averaged drag & setup keys \code{minimum\_cp}, \code{solver\_stabilization}, \code{unsteady\_pressure\_and\_kutta} and \code{reynolds\_averaged\_drag} & 26.120, 26.122, 26.123 \\
Advanced settings: wake relaxation, decay and agglomeration, and jet wakes & setup keys \code{wake\_relaxation}, \code{wake\_numerical\_relaxation}, \code{wake\_decay\_constant\_per\_m}, \code{wake\_streamwise\_agglomeration}, \code{jet\_wake\_decay\_normalized\_length} and \code{jet\_wake\_filaments\_grid\_induction} & none \\
Advanced settings: where the wake ends & setup key \code{wake\_termination\_revolutions} or \code{wake\_termination\_steps} & none \\
\bottomrule
\end{tabularx}}
\guisource{Flight conditions and solver}
\end{frame}

\begin{frame}{Flight conditions and solver: what a key takes (5/6)}
\relax
{\ftstablesize\renewcommand{\arraystretch}{1.05}
\begin{tabularx}{\linewidth}{@{}JKZ@{}}
\guihead
Advanced settings: the lift and separation models & setup keys \code{kutta\_joukowski\_lift}, \code{vorticity\_lift\_model}, \code{laminar\_separation}, \code{adverse\_gradient\_boundary\_layer}, \code{vortex\_ring\_normalization} and \code{axial\_separation\_families} & none \\
Advanced settings: rotor-induced velocities, field-grid refinement, the aeroelastic interpolation & setup keys \code{print\_rotor\_induced\_velocities}, \code{rotor\_induced\_velocity\_blending}, \code{adaptive\_field\_grid\_refinement} and \code{aeroelastic\_rbf\_type} & none \\
Create or delete a separation model: bulk, airfoil, axial vortex, cylindrical, Stratford & setup keys \code{bulk\_separation}, \code{airfoil\_separation}, \code{axial\_vortex\_separation}, \code{cylindrical\_bulk\_separation}, \code{stratford\_bulk\_separation} and \code{delete\_separations} & none \\
Set the Valarezo and crossflow separation criteria & setup keys \code{valarezo\_criterion}, \code{valarezo\_separation\_boundaries}, \code{crossflow\_separation\_boundaries}, \code{crossflow\_separation\_diameter}, \code{crossflow\_separation\_mean\_diameter} and \code{crossflow\_separation\_axisymmetric} & none \\
\bottomrule
\end{tabularx}}
\guisource{Flight conditions and solver}
\end{frame}

\begin{frame}{Flight conditions and solver: what a key takes (6/6)}
\relax
{\ftstablesize\renewcommand{\arraystretch}{1.05}
\begin{tabularx}{\linewidth}{@{}JKZ@{}}
\guihead
Mark thin or viscous-excluded boundaries, or set a surface roughness & setup keys \code{thin\_boundaries}, \code{viscous\_excluded} and \code{surface\_roughness} & 26.121 to 26.124, except two \\
Couple the blade's structure to the aerodynamics (FSI) & row key \code{FSI}, naming a structural configuration of \code{inputs/fsi/}, and fourteen row keys \code{FSI\_\textless{}PROPERTY\textgreater{}\_FACTOR}, each scaling one structural property; allowed on \code{steady}, \code{qsteady\_rotor} sector and \code{unsteady} rows, refused on \code{unsteady\_rotor} rows (Guide 8) & none \\
\bottomrule
\end{tabularx}}
\par\vspace{1mm}
\begin{alertblock}{The custom field is the flow's direction}
\small On 26.124 the angle of attack does not turn a custom field [2]: incidence goes into \code{vy} and \code{vz}, and a nonzero \code{ALPHA} or \code{BETA} beside \code{FREESTREAM} is refused at plan.
\end{alertblock}
\guisource{Flight conditions and solver}
\end{frame}
